% !PRINTED 320
\weberpage{320}
If one applies the substitution $\sigma$ to the indices $\xi, \eta$ of the $X$,
the indices of the $\Psi$ undergo [as in \S.\,80, (2)] the substitution
$\varrho^{-1}$, where $\varrho$ is the transposed substitution of $\sigma$.

But, corresponding to the substitutions (4), we have
\weq{6}{\varrho_1^{-1} = \begin{pmatrix}0, & 1\\1, & 0\end{pmatrix}, \quad
\varrho_2^{-1} = \begin{pmatrix}1, & 0\\1, & 1\end{pmatrix}.}

Assuming now, for the sake of simplicity, that the three quantities (5) are
different from zero, we set
\weq{7}{\begin{aligned}
v_{1,0} &= \frac{\Psi_{0,1}\Psi_{1,1}}{\Psi_{1,0}}
= \frac{\Psi_{0,1}\Psi_{1,0}\Psi_{1,1}}{\Psi_{1,0}^2};\\
v_{0,1} &= \frac{\Psi_{1,0}\Psi_{1,1}}{\Psi_{0,1}}
= \frac{\Psi_{0,1}\Psi_{1,0}\Psi_{1,1}}{\Psi_{0,1}^2};\\
v_{1,1} &= \frac{\Psi_{1,0}\Psi_{0,1}}{\Psi_{1,1}}
= \frac{\Psi_{1,0}\Psi_{0,1}\Psi_{1,1}}{\Psi_{1,1}^2}.
\end{aligned}}

The quantities $v_{\xi,\eta}$ remain unchanged under the substitutions of $Q$,
and undergo, through $\sigma_1, \sigma_2$, the permutations determined by (6)
\[
\begin{array}{l l l l}
 & v_{1,0}, & v_{0,1}, & v_{1,1}\\
\sigma_1 & v_{0,1}, & v_{1,0}, & v_{1,1}\\
\sigma_2 & v_{1,1}, & v_{0,1}, & v_{1,0}.
\end{array}
\]

From this it follows that the symmetric functions of the three quantities
$v_{1,0}, v_{0,1}, v_{1,1}$ remain unchanged under the permutations $P$, and
are consequently contained in the domain of rationality of the symmetric
functions of the $X$. The three quantities $v$ are therefore the roots of a
cubic equation, and there would be no difficulty in forming this cubic
resolvent of the general biquadratic equation.

But from (7) it follows further that
\[
v_{1,0}\,v_{0,1} = \psi_{1,1}^2, \quad v_{1,0}\,v_{1,1} = \psi_{0,1}^2,
\quad v_{0,1}\,v_{1,1} = \psi_{1,0}^2, % [sic] printed with lowercase psi, elsewhere capital Psi is used for these quantities
\]
and accordingly, if we further set the rational quantity
\[
X_{0,0} + X_{1,0} + X_{0,1} + X_{1,1} = a
\]
it follows from (5) that
\weq{8}{4X_{0,0} = a + \sqrt{v_{1,0}}\,\sqrt{v_{0,1}}
+ \sqrt{v_{1,0}}\,\sqrt{v_{1,1}} + \sqrt{v_{0,1}}\,\sqrt{v_{1,1}}.}

The expression (8) gives only four different values when we determine the
signs of the three square roots occurring in it in all possible ways, and one
can also con-
%CONT>>
